% Fragmento público generado desde propietarios íntegros.
% Residencia: 52.9.1.
% Fuente: ../incoming/radion_monografia_20260827/manuscrito/sections/10_minkowski_hodge_lorentz.tex:56-100
\subsubsection{Hodge Duality and Quarter-turn Operator}

With the signature \((-+++ )\) and an orientation fixed, the Hodge star on
two-forms satisfies
\begin{equation}
\star:\Lambda^2Q_4^*\to\Lambda^2Q_4^*,
\qquad \star^2=-I.
\label{eq:hodge-star}
\end{equation}
Thus, the quarter-turn expressed in the phase plane by
\(J^2=-I\) reappears in field space as a complex structure on
\(\Lambda^2\). The domains differ, but the genealogy of orientation
is preserved by the functor.

Let \(A_{\mathrm{em}}\) be a potential and
\begin{equation}
F=\diff A_{\mathrm{em}}.
\label{eq:F-dA}
\end{equation}
The decomposition of \(F\) relative to a temporal observer produces electric
and magnetic fields; \(\star F\) exchanges their roles with the sign fixed
by the orientation and signature. The constitutive polarity of
\eqref{eq:EM-conjugation} thus acquires a differential realization.

\paragraph{Electromagnetic and causal transduction chain}

The transport that must remain visible is
\begin{equation}
\Gamma_9
\longrightarrow\text{memory helix}
\longrightarrow(Q_4,B_{\mathrm{caus}},\star)
\longrightarrow A_{\mathrm{em}}
\longrightarrow F=\diff A_{\mathrm{em}}
\longrightarrow
m\nabla_u u=q(\iota_uF)^\sharp.
\label{eq:lorentz-chain}
\end{equation}
Each arrow adds structure: the holonomy preserves return; the quotient fixes
causality; Hodge fixes duality; the potential fixes the field; coupling
to charge produces dynamics.

The final equation is the Lorentz force in geometric form. The TPK helix
is its memory antecedent, not a physical orbit already moved by a field.
Dynamics appears only upon introducing \(F\), \(q\), and the connection \(\nabla\).

